% =============================================================================
% UNL Dissertation Template
% Based on nuthesis.cls (Hummel & Dieter, 2013)
% Author: Joe Sandow, University of Nebraska-Lincoln
% Chapter 2: Do Guns Make Us Safe?
% =============================================================================

\documentclass{nuthesis}

% -----------------------------------------------------------------------------
% Required packages
% -----------------------------------------------------------------------------

% Math
\usepackage{amsmath}
\usepackage{amsfonts}
\usepackage{bbm}                    % for \mathbbm{1} indicator function

% Font
\usepackage[sc,osf]{mathpazo}       % Palatino font, matches nuthesis example
\usepackage{microtype}              % Better typography

% Citations
\usepackage{natbib}

% Tables
\usepackage{booktabs}
\usepackage{threeparttable}         % For table notes
\usepackage{longtable}
\usepackage{tabularx}
\usepackage{multirow}
\usepackage{makecell}

% Figures
\usepackage{graphicx}
\usepackage{caption}
\usepackage{subcaption}
\usepackage{float}

% Layout
\usepackage{pdflscape}
\usepackage{rotating}
\usepackage{indentfirst}
\usepackage{paralist}

% TikZ (for diagrams if needed)
\usepackage{tikz}
\usetikzlibrary{trees,shapes,arrows,positioning}

% Fixes for memoir + hyperref compatibility
\usepackage{memhfixc}

% Hyperlinks (load last)
\usepackage{hyperref}
\hypersetup{
    colorlinks = true,
    linkcolor  = black,
    citecolor  = black,
    urlcolor   = black
}

\setlength{\footskip}{12pt}

% =============================================================================
% Document metadata — UPDATE THESE FIELDS
% =============================================================================

\begin{document}

\frontmatter

\title{Do Guns Make Us Safe? Causal Evidence from a County-Level Panel}
\author{Joe Sandow}
\adviser{Professor [Edward Balistreri]}          % UPDATE
\adviserAbstract{[Edward Balistreri]}            % UPDATE (first and last name only)
\major{Economics}
\degreemonth{[December]}                       % UPDATE e.g. May
\degreeyear{[2026]}                         % UPDATE e.g. 2026

% =============================================================================
% Title page (auto-generated by nuthesis)
% =============================================================================

\maketitle

% =============================================================================
% Abstract (max 350 words including title, name, etc.)
% =============================================================================

\begin{abstract}

This study estimates the causal effect of household firearm ownership rates 
on total firearm mortality using a county-level panel dataset spanning 
1999--2016. To address endogeneity and reverse causality, household firearm 
ownership rates are instrumented with the population-weighted leave-one-out 
mean murder arrest rate of neighboring counties. The preferred two-way fixed 
effects specification indicates that a one percentage point increase in the 
household firearm ownership rate causally increases the firearm death rate by 
approximately 0.065 additional deaths per 100,000 population --- an increase 
of roughly 0.85\% relative to the sample mean. These results are robust to 
weak-instrument concerns under Anderson-Rubin inference ($F = 56.68$, 
$p < 0.001$). Unlike much of the existing literature, which examines firearm 
ownership and homicide or suicide in isolation, this study examines total 
firearm mortality as the primary outcome, representing among the first 
IV-based estimates of the causal effect of ownership on total firearm 
mortality at the county level. Evidence on universal background check and 
concealed carry permit requirements is statistically inconclusive once 
standard errors are clustered at the state level, consistent with limited 
within-state policy variation over the sample period.

\end{abstract}

% =============================================================================
% Copyright page (optional)
% =============================================================================

\begin{copyrightpage}
% Auto-generates: Copyright © [year], [author]
% Leave blank or add text below if needed
\end{copyrightpage}

% =============================================================================
% Dedication (optional — delete environment if not needed)
% =============================================================================

\begin{dedication}
% e.g. To my family.
\end{dedication}

% =============================================================================
% Acknowledgments (optional — delete environment if not needed)
% =============================================================================

\begin{acknowledgments}
% Thank your adviser, committee, funding sources, etc.
\end{acknowledgments}

% =============================================================================
% Front matter tables
% =============================================================================

\tableofcontents
\listoffigures
\listoftables

% =============================================================================
% Main body
% =============================================================================

\mainmatter

% -----------------------------------------------------------------------------
% Chapter 2: Do Guns Make Us Safe?
% (If this is a standalone chapter submission use \chapter{}
%  If submitting the full dissertation this becomes \chapter{} 
%  preceded by any earlier chapters)
% -----------------------------------------------------------------------------

\chapter{Do Guns Make Us Safe? Causal Evidence from a County-Level Panel}
\label{chap:guns}

% -----------------------------------------------------------------------------
\section{Introduction}
\label{sec:intro}
% -----------------------------------------------------------------------------

Mass shooting events are often invoked in shorthand—Columbine, Sandy Hook, Uvalde, Pulse, Aurora, Parkland, Las Vegas—forming a macabre litany that blurs with other national tragedies, including Pearl Harbor, Oklahoma City, and September 11. Responses to the assertion that ``guns make us safe, and therefore more guns make us safer'' may reveal more about lived experiences than about the claim's empirical validity. For those who live in communities where law enforcement or emergency services reliably arrive within three to four minutes, the necessity of firearm ownership may seem unwarranted. Conversely, in communities where emergency response times range from four to eighteen minutes—whether due to geography, policy, or resource allocation—the benefits and risks of firearm ownership may be evaluated differently. To adjudicate these competing intuitions, credible causal evidence is required. This study provides such evidence. 
This paper estimates the causal effect of firearm prevalence on total firearm mortality. Using a county-level panel dataset spanning 1999–2016, the analysis employs an instrumental variable strategy to address endogeneity and reverse causality in the relationship between firearm ownership and firearm mortality. Unlike much of the existing literature, which disaggregates firearm deaths into suicides, homicides, or accidental incidents to address outcome heterogeneity, this study examines total firearm mortality as the primary outcome. The preferred two-way fixed effects specification indicates that a one percentage point increase in the household firearm ownership rate causally increases the firearm death rate by approximately 0.065 additional deaths per 100,000 population — an increase of roughly $0.85\%$ relative to the sample mean. These findings contribute to a literature that has largely examined firearm ownership and homicide or suicide in isolation, and represent among the first IV-based estimates of the causal effect of ownership on total firearm mortality at the county level.


% -----------------------------------------------------------------------------
\section{Literature Review}
\label{sec:lit}
% -----------------------------------------------------------------------------

 The literature on firearms is broadly divided into two primary traditions: criminology and public health, with contributions from economists and sociologists aligning with one of these approaches. The majority of existing studies rely on correlational designs and therefore cannot fully address concerns related to endogeneity and reverse causality. Much of the early research emerged from criminology reflecting the persistently high rates of firearm homicide in the United States \citep{kleck2017point, kleck1993impact}.
 
Consistent with this focus, numerous studies document a strong positive correlation between firearm ownership and homicide rates. Some estimates approach a one-to-one relationship, with a one percent increase in firearm ownership associated with roughly a 0.9 percent increase in homicide rates \citep{siegel2013relationship}. Accordingly, many studies report positive associations between firearm ownership and a range of crime outcomes \citep{cook2006social, cook1979effect, duggan2001more, hemenway2000firearm, kaplan1998firearm, miller2002rates, miller2007state, sorenson2001handgun, koenig2016dynamics}. A common empirical strategy within this literature is to use the proportion of firearm suicides as a proxy for household firearm ownership, given the absence of comprehensive administrative ownership data. The causal interpretation of these associations is contested, however. Employing a LATE framework with county-level cross-sectional data and instrumenting gun prevalence with outdoor sports magazine subscriptions, presidential voting preferences, and military veteran shares, \citet{kovandzic2013estimating} find no significant positive effect of gun prevalence on homicide rates, interpreting their results as evidence that the deterrent effect of noncriminal gun ownership offsets any violence-facilitating effect.

Beyond homicide outcomes, some studies estimate the frequency with which firearms are used defensively to prevent crime. These estimates rely almost exclusively on self-reported survey data, as no comprehensive administrative measure of defensive firearm use currently exists. Approximately 30\% of firearm owners report having used a firearm defensively to prevent injury to themselves or a family member, a figure that has remained relatively stable for nearly three decades \citep{english20212021, kleck1993impact}. Evidence regarding concealed carry laws is mixed. Some studies find associations between concealed carry laws and reductions in certain crime categories \citep{lott2013more}\footnote{Lott's empirical methodology and data have been subject to substantial scrutiny in the literature.}, while others report inconclusive or mixed findings \citep{koenig2016dynamics}. More recent evidence suggests that concealed carry permits have no statistically significant effect on gun-assault hospital visits, though increases in permit issuance are associated with higher rates of gun-related hospital admissions \citep{johnson2023armed}.

In contrast to the criminology tradition, research in public health has focused primarily on suicide and injury. Research examining firearms from a public health perspective has frequently centered on suicide (Figure \ref{fig: self_harm} ) In some jurisdictions, firearm suicides exceed firearm homicides by a factor of ten or more (Figure \ref{fig:sh_homicide_state}). Areas with higher firearm ownership rates exhibit higher rates of firearm suicide attempts, and suicide attempts involving firearms are substantially more likely to result in death \citep{briggs2014firearms}. Recent work using autoregressive cross-lagged models provides evidence on the temporal ordering of this relationship: increases in household firearm ownership rates precede increases in firearm suicide rates, while the reverse — firearm suicide rates predicting subsequent changes in ownership — is weaker in magnitude \citep{morral2024household}. This pattern is consistent with a causal interpretation in which firearm availability elevates suicide risk rather than elevated suicide risk driving ownership. Parallel to the suicide literature, researchers have examined firearm exposure and injury among children.

In April 2022, the \emph{New England Journal of Medicine} (NEJM) published an article reporting that firearm injuries surpassed motor vehicle accidents as the leading cause of death among children in 2017 \citep{doi:10.1056/NEJMp2200169}. Less than a month later, NEJM published a second article comparing firearm-related deaths to multiple leading causes of mortality, concluding that firearm deaths did not surpass motor vehicle accidents until 2019 \citep{doi:10.1056/NEJMc2201761}. Figure \ref{fig:child} replicates this comparison using CDC WONDER data, confirming the crossover in approximately 2019--2020. Although the precise timing can be debated, the broader conclusion is consistent: firearm-related injuries have been the leading cause of death among individuals younger than 19 since at least 2020. These studies contributed to the expanding literature on firearms and child mortality. Earlier research has examined several related questions, including whether Child Access Prevention (CAP) laws — designed to hold parents legally accountable for secure firearm storage — are effective \citep{desimone2013child}. \citet{desimone2013child} document a reduction in nonfatal firearm-related injuries among children under age 18 associated with CAP law adoption, using hospital discharge data — an outcome measure distinct from the mortality-focused analyses that dominate this literature. More broadly, \citet{tashiro2016effect} examine the relationship between state-level firearm regulations and hospital admissions for firearm-related injuries among children, finding that restrictive and permissive states accounted for comparable shares of total injuries — 47\% and 48\% respectively — suggesting limited evidence of a differential effect of general gun control measures on pediatric firearm injuries.

Other studies have examined firearm-related injuries treated at hospitals and attempted to quantify the associated economic costs \citep{de2020association, jehan2018burden}. Both studies rely on single-year cross-sectional data — \citet{jehan2018burden} using 2011 and \citet{de2020association} using 2014 — which precludes identification of causal effects from within-state policy variation over time. Both find that states with more restrictive firearm legislation are associated with lower rates of firearm-related injuries and reduced associated medical costs, though the cross-sectional designs cannot rule out confounding from pre-existing state-level differences.

Despite this extensive literature, relatively few studies examine the relationship between overall firearm ownership rates and total firearm mortality within a unified causal framework. One recent contribution analyzes concealed carry permit laws in Pennsylvania using several instrumental variable strategies. Drawing on evidence that local crime risk causally influences firearm-related decisions \citep{depew2019decision}, \citet{johnson2023armed} instrument for concealed carry permit rates using neighboring county murder rates and find that permit laws are associated with reductions in firearm-related injuries. However, that analysis is geographically limited to a single state and covers the period 2012--2020. Building on this approach, the present study employs neighboring county murder arrest rates as an instrument for firearm prevalence. This strategy enables identification using all fifty states over the period 1999--2016, substantially expanding both the geographic scope and time horizon relative to Johnson et al. The remainder of the paper proceeds as follows. Section 2 describes the data. Section 3 outlines the identification strategy. Section 4 presents the results, and Section 5 concludes.

% Figures referenced in literature review:

\begin{figure}[htbp]
    \centering
    \includegraphics[width=\textwidth]{idaho_mon_suicides_homicides.pdf}
    \caption{Firearm suicides and firearm homicides in Idaho and Montana, 
    1999--2020. In both states, firearm suicides exceed firearm homicides 
    by a factor of ten or more in most years. Similar patterns are observed 
    in Vermont, North Dakota, Wyoming, New Hampshire, and South Dakota, 
    where firearm homicides are suppressed or near zero across the entire 
    sample period.}
    \label{fig:sh_homicide_state}
\end{figure}

\begin{figure}[htbp]
    \centering
    \includegraphics[width=\textwidth]{self_harm_v_other.pdf}
    \caption{All Other includes firearm deaths excluding self-harm, 
    including accidental deaths and deaths resulting from legal intervention.}
    \label{fig:self_harm}
\end{figure}

\begin{figure}[htbp]
    \centering
    \includegraphics[width=\textwidth]{NEJM_reproduction.pdf}
    \caption{Leading causes of death among children and adolescents in the 
    United States, 1999--2020, expressed as deaths per 100,000. Data are 
    drawn from the CDC WONDER Underlying Cause of Death database. This 
    figure replicates the structure of Figure 1 in 
    \citet{doi:10.1056/NEJMc2201761}, extending their analysis using the 
    same underlying data source.}
    \label{fig:child}
\end{figure}

% -----------------------------------------------------------------------------
\section{Data}
\label{sec:data}
% -----------------------------------------------------------------------------

This study uses firearm mortality data from the CDC WONDER Underlying Cause of Death database for the period 1999–2020, though the estimation sample is restricted to 1999–2016 by the availability of firearm ownership estimates. These data are merged with state-level household firearm ownership estimates compiled by the RAND Corporation \citep{TL-354-LJAF}. County adjacency is determined using the 2010 U.S. Census county adjacency file, which relies on Federal Information Processing Standards (FIPS) county codes. County-level arrest data are obtained from the Inter-university Consortium for Political and Social Research (ICPSR). ICPSR compiles these data using the FBI's Uniform Crime Reporting (UCR) system, which provides county-level crime statistics through 2016, excluding 2015. For 2015 and, in the descriptive analysis, 2017--2020, arrest counts are constructed by aggregating agency-level UCR reports to the county level using the Law Enforcement Agency Identifiers Crosswalk (LEAIC). Because Florida does not report complete crime statistics to the FBI, county-level Total Index Crime data for Florida are obtained directly from the Florida Department of Law Enforcement (1999--2020).

Firearm mortality data are obtained from the CDC WONDER Underlying Cause of Death database for the years 1999–2020. The database is compiled from official death certificates, providing a high degree of administrative accuracy. However, two limitations are particularly relevant for the present analysis. First, mortality coding relies on the primary cause of death listed on the certificate. In many states, only a single underlying cause is recorded. Consequently, if an individual sustains a firearm injury but subsequently dies from a secondary medical complication (e.g., cardiac arrest), the death may be classified under the proximate medical cause rather than the precipitating firearm injury. Although such cases are likely infrequent, this coding structure may lead to some degree of undercounting of firearm-related mortality. Second, the CDC suppresses county-year observations in which fewer than ten deaths are recorded in order to protect individual confidentiality \citep{cdc_wonder_suppression}.

Although the CDC WONDER database allows highly granular queries (e.g., county-by-month), most such observations are suppressed due to confidentiality restrictions. To mitigate suppression, the analysis is conducted at the county-year level. As a result, demographic disaggregation, such as by sex assigned at birth or age, is not feasible at the county-year level. National annual death counts by firearm-related ICD-10 cause-of-death code were obtained for 1999–2020 (Figure \ref{fig: self_harm}). At the county level, only 47\% of counties have at least one non-suppressed county-year observation over the sample period. Consequently, at least 53\% of counties are fully suppressed across all 22 years. The balanced panel would contain 3,147 counties observed over 22 years, yielding 69,234 potential county-year observations. Due to suppression, the estimating sample includes approximately 944 counties — roughly 30\% of the balanced panel — in any given year. Because suppression disproportionately affects low-population areas, the resulting sample overrepresents more populous, urban counties. Despite limited county coverage, the observed counties account for approximately 75\% of the U.S. population in 2000 and approximately 82\% by the end of the sample period, representing upper bounds on annual population coverage (Figure \ref{fig:pop_cov}). On average, the estimating sample accounts for approximately 79.5\% of the U.S. population annually.

A central empirical challenge in this literature is the absence of reliable administrative data on firearm stocks. The total number of firearms currently in private circulation in the United States is unknown. Frequently cited estimates suggest that approximately 400 million firearms are privately owned in the United States — more firearms than residents. This estimate is typically constructed by aggregating domestic production, subtracting net exports, and applying an assumed depreciation rate to account for firearms exiting circulation. However, empirical evidence on firearm depreciation rates is limited \citep{alvazzi2020small}. Consequently, aggregate firearm stock estimates remain subject to substantial measurement uncertainty.

The most comprehensive estimates of household firearm ownership are provided by the RAND Corporation, which compiled 22 distinct measures of firearm ownership spanning 1980--2016. These measures incorporate both self-reported survey data and structural proxies, including permit issuance and subscriptions to firearms-related publications, to estimate the proportion of households owning firearms at the state level. Among the self-reported inputs are three Behavioral Risk Factor Surveillance System (BRFSS) waves that included firearm ownership questions (2001, 2002, and 2004). The RAND estimates are constructed at the state level, and no comparable, reliable county-level ownership measures currently exist \citep{TL-354-LJAF}. Accordingly, each county is assigned the household firearm ownership rate corresponding to its state. The implications of this assignment for measurement error and identification are discussed below.

Although firearm ownership rates likely vary between rural and urban counties, the use of state-level ownership estimates for county-level analysis can be justified both algebraically and empirically. Suppression in the mortality data disproportionately affects sparsely populated counties. For example, Vermont—the most rural state, with 64.9\% of its population classified as rural—has approximately 85\% of its county-year observations suppressed. In contrast, states with the largest rural populations in absolute terms (e.g., Texas, North Carolina, Pennsylvania, and Ohio) have substantially lower rural population shares: 15\% in Texas, 32\% in North Carolina, and approximately 23\% in Pennsylvania and Ohio. Consider North Carolina as an illustrative case. Over the period 1999–2016, the mean household firearm ownership rate (HFR) is 37\%. Let $r$ denote the rural ownership rate, $\mu$ the urban ownership rate, and $\alpha = 0.32$ the rural population share. The state mean satisfies
\begin{align*}
    0.37 &= \alpha r + (1-\alpha)\mu
\end{align*}
Even under an extreme upper bound in which $r=1$, the implied urban ownership rate is 
$\mu=0.0641$. This bound illustrates that the state-level mean places substantial weight on the urban ownership rate. Because the estimating sample overrepresents urban counties due to suppression, the state-level ownership measure is likely closer to the ownership rates prevailing in the observed (primarily urban) counties than to those in sparsely populated rural areas. To the extent that within-state heterogeneity remains, the resulting measurement error is unlikely to be systematically correlated with the instrument and would therefore be expected to attenuate estimated effects.

Household firearm ownership rates exhibit a general downward trend from 1999 to 2010 (Figure \ref{fig:HFR}), continuing a broader decline that began in the 1980s. In 2011, ownership rates increased modestly, followed by a more pronounced rise in 2012, from approximately 0.39 to 0.41. This increase was followed by three consecutive years of decline, reaching approximately 0.37, before rising again to 0.40 in 2016. These fluctuations coincide with similar movements in the national firearm mortality rate during the same period.

% Figures referenced in data section:

\begin{figure}[htbp]
    \centering
    \includegraphics[width=\textwidth]{percent_covered.pdf}
    \caption{Due to suppression in the CDC WONDER data, the share of the 
    U.S. population represented in the estimating sample varies by year. 
    The reported values represent upper bounds on annual population 
    coverage. The upward trend is consistent with increasing population 
    concentration in urban counties over time.}
    \label{fig:pop_cov}
\end{figure}

\begin{figure}[htbp]
    \centering
    \includegraphics[width=\textwidth]{HFR.pdf}
    \caption{Household firearm ownership rates exhibit a continued decline 
    from 1999 through 2010, extending the broader downward trend that began 
    in 1980. Ownership rises in 2011 and 2012 before declining again 
    through 2016.}
    \label{fig:HFR}
\end{figure}

\begin{figure}[htbp]
    \centering
    \includegraphics[width=\textwidth]{mean_deaths_per_hund.pdf}
    \caption{Mean firearm deaths per 100,000 population by year, 1999--2020. 
    Estimates are derived from state-level data and reflect the national 
    unweighted mean across states.}
    \label{fig:death_rate}
\end{figure}

% -----------------------------------------------------------------------------
\section{Identification}
\label{sec:id}
% -----------------------------------------------------------------------------

Empirical analyses typically begin with a formal specification of the econometric model. However, because the baseline specification is straightforward, the discussion begins with the identification strategy and the construction of the instrumental variable. The need for an instrumental variable arises from concerns about reverse causality. For example, if firearm violence increases in a neighborhood, residents may respond by purchasing firearms for protection, generating endogeneity in observed ownership rates. An ideal instrument would affect household firearm ownership while having no direct effect on firearm mortality. In studies focused narrowly on gun crime, researchers often use firearm suicide rates as an instrument for firearm ownership. However, because the present analysis examines total firearm mortality, including suicides, this approach would violate the exclusion restriction and reintroduce endogeneity concerns. A substantial literature documents a correlation between firearm ownership rates and crime rates \citep{cook2006social, cook1979effect, duggan2001more, hemenway2000firearm, kaplan1998firearm, miller2002rates, miller2007state, sorenson2001handgun, koenig2016dynamics}. \citet{depew2019decision} provide causal evidence that local homicide rates influence individual decisions to apply for concealed carry permits, establishing a link between neighborhood crime and firearm-related behavior. Building on this relationship, \citet{johnson2023armed} construct a leave-one-out instrument using neighboring county murder rates to identify the effect of concealed carry permit rates on firearm-related injuries. The present study adopts a similar strategy but uses murder arrest rates rather than murder rates, as arrest data are available at the county level through the UCR system, and instruments for household firearm ownership rates rather than concealed carry permit rates, thereby broadening both the outcome of interest and the geographic scope of identification.

A natural concern is that the proposed instrument may violate the exclusion restriction. For example, an individual residing in a suburban county may commit a homicide in a neighboring urban county but be arrested in their county of residence. While such cross-county incidents may occur, several conditions must hold for this mechanism to generate systematic bias. First, the offender must reside in a county contiguous to the one in which the crime occurred. Second, such cross-county arrest patterns would need to be sufficiently prevalent and geographically systematic to bias the instrument. Empirical evidence suggests that most arrests occur within the offender's county of residence, indicating that criminal activity and arrest outcomes are predominantly local phenomena \citep{johnson2023armed}.

The exclusion restriction could also fail if criminal activity systematically occurs across county borders, such that offenses are committed in one county but arrests consistently occur in another. However, available evidence suggests that most violent crime is geographically localized rather than systematically cross-jurisdictional. While organized criminal networks may operate across jurisdictions, such activity would need to be sufficiently widespread and systematically aligned with county borders to generate meaningful bias in the instrument. Empirical evidence indicates that arrests overwhelmingly occur within the county in which the offense was committed, limiting the scope for systematic cross-county misalignment \citep{johnson2023armed}. Accordingly, the instrument is constructed as the population-weighted leave-one-out mean murder arrest rate among counties neighboring county ii
i, as defined formally below.
    \begin{align}
        WAR_{it}=\sum_{j\neq i}^{n}p_{jt}ar_{jt}
    \end{align}

\noindent where $ar_{jt}$ denotes the murder arrest rate in county $j$ at time $t$, and $p_{jt}$ denotes the population weight of county $j$ among all counties adjacent to $i$.  

Returning to Figure \ref{fig:death_rate}, the time series warrants closer examination. Through 2010, the firearm death rate per 100,000 remained relatively stable at approximately 11. Beginning in 2010, the firearm death rate began a sustained upward trend that continued through 2020. After peaking in 2020, the firearm death rate declined, though not to pre-2010 levels.
\begin{quote}
    Private citizens have the right under the Second Amendment to possess an ordinary type of weapon and use it for lawful, historically established situations such as self-defense in a home, even when there is no relationship to a local militia.%
    \footnote{\textit{District of Columbia v. Heller}, 554 U.S. 570 (2008).}
\end{quote}
The decision recognized an individual right to possess firearms for self-defense, though its immediate application was limited to federal jurisdictions. It was not until 2010 that the Supreme Court found:
\begin{quote}
    The Due Process Clause of the Fourteenth Amendment extends the Second Amendment's right to keep and bear arms to the states, at least for traditional, lawful purposes such as self-defense.%
    \footnote{\textit{McDonald v. City of Chicago}, 561 U.S. 742 (2010).}
\end{quote}
The Court’s decision in McDonald incorporated the Second Amendment against the states, thereby invalidating certain state-level firearm restrictions. Figure \ref{fig: HFR} indicates increases in estimated household firearm ownership in 2011 and 2012, coinciding with the initial rise in firearm mortality shown in Figure \ref{fig:death_rate}. However, attributing these changes to the Court’s decisions would require ruling out alternative explanations for contemporaneous increases in firearm ownership and mortality. Such a causal interpretation is further complicated by evidence that suicides account for the majority of the increase in firearm mortality during 2011–2012 (Figure \ref{fig: self_harm}). The period also coincides with broader social changes, including the rapid expansion of social media platforms. Prior research has linked increased social media use to elevated suicide risk, particularly among adolescent females \citep{jaycox2024social, massing2020cybervictimization}. Figure \ref{fig:sh_gender} shows a modest increase in firearm suicides among individuals assigned female at birth beginning around 2010. The corresponding increase among individuals assigned male at birth is substantially larger.

In the panel setting, inference must account for serial correlation within states and common shocks over time. The choice of functional form also matters: a log specification is often convenient for rate outcomes, but it is undefined when the outcome is zero. In this application, some county-year observations record zero firearm deaths, implying a zero firearm death rate. These zeros reflect valid realizations of the outcome and should be retained in the estimating sample. Accordingly, the baseline specification is estimated in levels rather than log–log form.

\begin{align}
    d\_rate_{i,t} &= \beta_{1}HFR_{s,t} + \gamma_{s} + \delta_{t} +\mathbbm{1}\{universl_{s,t}\} + \mathbbm{1}\{permit_{s,t}\} +\epsilon_{i,t}
\end{align}

\noindent where $d\_rate_{i,t}$ denotes firearm deaths per capita in county $i$ at time $t$; $HFR_{s,t}$ denotes the state-level household firearm ownership rate; $s$ indexes states; and $t$ indexes years. $\mathbbm{1}\{universl_{s,t}\}$ and $\mathbbm{1}\{permit_{s,t}\}$ are indicator variables equal to one if state $s$ has universal background checks or permit requirements, respectively, in year $t$, and zero otherwise. Over the period 1999–-2016 there is limited within-state variation in either indicator. During this period, states with universal background checks typically also required permits, and such policies were rarely repealed. Consequently, these indicators primarily capture cross-state differences rather than within-state policy changes. $\gamma_{s}$ and $\delta_{t}$ denote state and year fixed effects, respectively.

\begin{align}
    HFR_{s,t} = \varphi_{1}WAR_{i,t} + \gamma_{s} + \delta_{t} 
    + \mathbbm{1}\{universal_{s,t}\} 
    + \mathbbm{1}\{permit_{s,t}\} 
    + \nu_{i,t}
\end{align}

\noindent $WAR_{i,t}$ denotes the population-weighted mean murder arrest rate among counties contiguous to county $i$, excluding county $i$ itself — the leave-one-out instrument. In the first-stage regression, the coefficient on $WAR_{i,t}$ is positive and statistically significant at the 5\% level with state-clustered standard errors ($\hat{\varphi}_{1} = 0.000182; se = 0.0000880$), consistent with the instrument shifting household firearm ownership rates in the expected direction. The first-stage F-statistic on the excluded instrument is $16.18$ ($p < 0.001$) without clustering, exceeding the conventional threshold of 10 \citep{stock2002testing}. When standard errors are clustered at the state level, however, the effective first-stage F-statistic falls to approximately 4.27, raising concerns about instrument strength under clustered inference. To address this, the Anderson-Rubin test \citep{anderson1949estimation} — which provides valid inference regardless of instrument strength — is employed as a robustness check. The Anderson-Rubin test rejects the null hypothesis of no treatment effect ($F = 56.68, p < 0.001$), with a 95\% confidence interval of $[0.004, 0.013]$, consistent with the 2SLS point estimate and supporting a positive causal effect of household firearm ownership on firearm mortality. Subsequent specifications include an additional indicator variable $\mathbbm{1}\{decision_{i,t}\}$, equal to one for observations after 2010 (and after 2008 for Washington, D.C.), to capture the structural break in firearm mortality coinciding with the \textit{McDonald v. City of Chicago} decision.\footnote{\textit{McDonald v. City of Chicago}, 561 U.S. 742 (2010).} In these specifications, $\mathbbm{1}\{decision_{i,t}\}$ replaces the year fixed effects while state fixed effects are retained.

Identifying a valid instrument for household firearm ownership is inherently difficult: any variable strongly correlated with ownership rates is likely also correlated with the structural disturbance $\epsilon_{i,t}$, violating the exclusion restriction. Hunting-related variables have appeared in prior work as proxies or instruments for firearm prevalence — including outdoor sports magazine subscriptions \citep{kovandzic2013estimating} and hunting license rates \citep{siegel2013relationship} — reflecting the strong correlation between hunting participation and household firearm ownership. For example, a safety-motivated purchaser is unlikely to hold a hunting license, illustrating that hunting-related variables may not capture the full range of ownership motivations relevant to the present context. More fundamentally, hunting participation is itself embedded in the broader cultural determinants of firearm prevalence and, potentially, firearm mortality — suggesting that hunting licenses may fail the exclusion restriction through channels independent of ownership rates.

An alternative instrument employed in prior work is the homicide rate in neighboring counties \citep{johnson2023armed}. This approach performs reasonably well in urban areas but faces limitations in rural counties, where homicide counts are low and CDC suppression is prevalent, reducing the precision and coverage of the instrument. The leave-one-out murder arrest rate instrument faces the same limitation in rural counties; however, because CDC suppression excludes most rural county-year observations from the estimating sample, this constraint has limited practical impact on the analysis. Neighboring county suicide rates are similarly unsuitable as an instrument, as they would violate the exclusion restriction through the firearm substitution effect. Firearm ownership increases the likelihood of firearm suicide among suicidal individuals \citep{briggs2014firearms}, and recent evidence suggests ownership causally precedes increases in firearm suicide rates \citep{morral2024household} — meaning that any instrument correlated with suicide rates would also be correlated with the outcome through this channel. The number of taxidermists per capita has also been proposed as an instrument for firearm ownership, given that hunters frequently use taxidermy services — establishing a link between taxidermist prevalence and hunting activity, and by extension, firearm ownership.

Although taxidermists represent an ancillary component of hunting culture, the channel through which taxidermist prevalence would affect firearm mortality, independent of ownership rates, is not apparent, suggesting the instrument may satisfy the exclusion restriction but raise concerns about relevance. A practical obstacle is that the North American Industry Classification System (NAICS) does not provide a dedicated code for taxidermists, instead classifying them under the broader category of independent artists, which encompasses actors, writers, and directors — making it difficult to isolate taxidermist counts from aggregate data. Whether this data limitation could be resolved through alternative sources, such as state-level business licensing records, remains an avenue for future research. Despite the instrument strength concerns discussed above, the leave-one-out neighboring county murder arrest rate remains attractive relative to available alternatives: it is derived from administrative data sources, its exclusion restriction is more plausibly defensible than hunting-related or suicide-based instruments, and the Anderson-Rubin test provides weak-instrument-robust evidence of a significant treatment effect.

% Figures referenced in identification section:

\begin{figure}[htbp]
    \centering
    \includegraphics[width=\textwidth]{SH_gender.pdf}
    \caption{Firearm suicides by sex assigned at birth, 1999--2020. The 
    peak number of firearm suicides among individuals assigned female at 
    birth occurred in 2017--2018 at approximately 3,500 deaths, 
    representing roughly 15\% of firearm suicides in those years.}
    \label{fig:sh_gender}
\end{figure}

% -----------------------------------------------------------------------------
\section{Results}
\label{sec:results}
% -----------------------------------------------------------------------------

Table \ref{tab:main_results} presents IV estimates of the effect of household firearm ownership on firearm mortality across three specifications. Column (1) is a pooled IV specification with no fixed effects. Column (2) adds state and year fixed effects and represents the preferred specification. Column (3) replaces year fixed effects with the post-2010 indicator $\mathbbm{1}\{decision_{i,t}\}$. The estimating sample consists of 19,922 county-year observations drawn from the 1999--2016 period, the latter bound determined by the availability of RAND household firearm ownership estimates. As established in the identification section, the first-stage coefficient on $WAR_{i,t}$ is positive and statistically significant ($\hat{\varphi_{1}} = 0.000182; se = 0.0000880$), and the Anderson-Rubin test confirms instrument validity under weak-instrument-robust inference ($F = 56.68, p < 0.001$). Turning to the second stage, the preferred specification in Column (2) indicates that a one percentage point increase in the household firearm ownership rate causally increases the firearm death rate by approximately $6.53 \times 10^{-5}$ deaths per capita — equivalent to roughly 0.065 additional firearm deaths per 100,000 population — after absorbing state and year fixed effects. Relative to the sample mean of 7.64 firearm deaths per 100,000 population, this corresponds to an increase of approximately $0.85\%$. The post-2010 indicator $\mathbbm{1}\{decision_{i,t}\}$ is positive and statistically significant in Column (3), consistent with the structural break in firearm mortality identified in the descriptive analysis. The policy indicators $\mathbbm{1}\{universal_{s,t}\}$ and $\mathbbm{1}\{permit_{s,t}\}$ are statistically insignificant across all specifications once standard errors are clustered at the state level, consistent with the limited within-state variation in these indicators over the 1999--2016 period documented in the data section.

The positive and statistically significant coefficient on $\mathbbm{1}\{decision_{i,t}\}$ in Column (3) is consistent with the structural break in firearm mortality identified in the descriptive analysis, in which the national firearm death rate remained stable at approximately 7.64 per 100,000 through 2010 before rising steadily through the end of the sample period. The timing of this break coincides with the Supreme Court's ruling in \textit{McDonald v. City of Chicago} (2010),\footnote{\textit{McDonald v. City of Chicago}, 561 U.S. 742 (2010).} which extended the individual right to firearm possession established in \textit{District of Columbia v. Heller} (2008)\footnote{\textit{District of Columbia v. Heller}, 554 U.S. 570 (2008).} to the states. However, as discussed in the conclusions, the precise mechanism driving the post-2010 trend is not fully resolved by the present analysis, and $\mathbbm{1}\{decision_{i,t}\}$ should be interpreted as capturing a broad structural shift in firearm mortality rather than the causal effect of a single legal ruling.

% Main results table placeholder:
\begin{table}[!htbp]
\centering
\caption{IV Estimates of the Effect of Household Firearm Ownership 
on Firearm Mortality}
\label{tab:main_results}
\begin{tabular}{lccc}
\toprule
 & Pooled IV & Two-Way FE & State FE + Decision \\
 & (1) & (2) & (3) \\
\midrule
HFR & $-$0.005 & 0.007$^{*}$ & 0.003$^{**}$ \\
    & (0.006)  & (0.004)     & (0.002) \\[0.5em]
$\mathbbm{1}\{universal\}$ & $-$0.0002 & 0.0001 & 0.00004 \\
    & (0.0004) & (0.0001) & (0.00003) \\[0.5em]
$\mathbbm{1}\{permit\}$ & $-$0.001 & 0.0002 & 0.0001 \\
    & (0.001)  & (0.0002) & (0.0001) \\[0.5em]
$\mathbbm{1}\{decision\}$ & & & 0.00002$^{**}$ \\
    & & & (0.00001) \\
\midrule
State FE           & No           & Yes          & Yes \\
Year FE            & No           & Yes          & No \\
SE Clustering      & State        & State        & State \\
Instrument         & $WAR_{it}$   & $WAR_{it}$   & $WAR_{it}$ \\
First-Stage F      & 60.09        & 16.08        & 38.28 \\
Effective F        & ---          & 5.14         & --- \\
Wu-Hausman $p$     & $<$0.001     & $<$0.001     & $<$0.001 \\
Sargan             & Exactly ID   & Exactly ID   & Exactly ID \\
Observations       & \multicolumn{3}{c}{19,922} \\
\bottomrule
\end{tabular}
\begin{minipage}{\textwidth}
\footnotesize
\vspace{4pt}
\textit{Notes:} All specifications instrument household firearm ownership 
rates ($HFR_{s,t}$) with the leave-one-out population-weighted mean murder 
arrest rate of neighboring counties ($WAR_{it}$). Column (1) is pooled IV 
with no fixed effects. Column (2) adds state and year fixed effects and 
represents the preferred specification. Column (3) replaces year fixed 
effects with $\mathbbm{1}\{decision_{i,t}\}$, equal to one for observations 
after 2010. $\mathbbm{1}\{universal_{s,t}\}$ and $\mathbbm{1}\{permit_{s,t}\}$ 
indicate states with universal background check and concealed carry permit 
requirements, respectively. The conventional first-stage F is reported 
without clustering; the Olea-Pflueger effective F accounts for the clustering 
structure and is reported for Column (2). The Wu-Hausman test rejects the 
null of exogeneity in all specifications. The model is exactly identified in 
all specifications and the Sargan overidentification test is not applicable. 
Standard errors clustered at the state level in parentheses. 
$^{*}p<0.1$, $^{**}p<0.05$, $^{***}p<0.01$.
\end{minipage}
\end{table}  
% OR paste table directly here



% -----------------------------------------------------------------------------
\section{Conclusion}
\label{sec:conclusion}
% -----------------------------------------------------------------------------

This study examines the causal relationship between household firearm ownership rates and total firearm mortality using a county-level panel dataset spanning 1999--2016. Employing a leave-one-out instrumental variable strategy — the population-weighted mean murder arrest rate of neighboring counties — and instrumenting for state-level household firearm ownership rates, the preferred two-way fixed effects specification indicates that a one percentage point increase in the household firearm ownership rate causally increases the firearm death rate by approximately $6.53 \times 10^{-5}$ deaths per capita, equivalent to roughly 0.065 additional firearm deaths per 100,000 population, or approximately $0.85\%$ relative to the sample mean. These results are robust to weak-instrument concerns, as the Anderson-Rubin test rejects the null hypothesis of no treatment effect ($F = 56.68$, $p < 0.001$) with a 95\% confidence interval of $[0.004, 0.013]$.

The evidence on universal background check and concealed carry permit requirements is less conclusive. These indicators are statistically insignificant once standard errors are clustered at the state level, consistent with the limited within-state policy variation over the 1999–-2016 period. This finding adds to the already contested empirical literature on the effectiveness of these two policy instruments \citep{kovandzic2013estimating, siegel2013relationship}. The present results neither confirm nor rule out meaningful effects of these policies — rather, they suggest that the identifying variation available in this sample is insufficient to precisely estimate their causal effects, and that panel designs with richer within-state policy variation would be better suited to that question.

The post-2010 structural break in firearm mortality, captured by $\mathbbm{1}\{decision_{i,t}\}$, is statistically significant and consistent with the expansion of individual firearm rights to the states following \textit{McDonald v. City of Chicago}.\footnote{\textit{McDonald v. City of Chicago}, 561 U.S. 742 (2010).} However, a closer examination of the descriptive evidence complicates a straightforward \textit{Heller} and \textit{McDonald} interpretation. Firearm suicides among individuals assigned male at birth began rising as early as 2006, four years before \textit{McDonald} and two years before \textit{Heller}, while all other firearm deaths were declining over the same period, with the net effect on the overall firearm death rate not materializing until 2010. The causal effect of firearm ownership on total mortality established by the IV estimates is therefore distinct from the question of what drove the post-2010 structural break — the former is identified through within-county variation in ownership rates instrumented by neighboring arrest rates, while the latter reflects aggregate time-series dynamics that the present specification captures but does not fully decompose.
The year 2006 predates the founding of Instagram and precedes Facebook's emergence as the dominant social media platform, which occurred around 2008. The bulk of veterans from the Second Iraq War began returning in 2007, the same year the Great Recession commenced. Each of these developments has been independently linked to elevated suicide risk in the existing literature, and their confluence around 2006–-2008 warrants further investigation as a potential explanation for the observed increase in male firearm suicides. Disentangling the relative contributions of expanded firearm access, social media exposure, veteran reintegration, and economic distress to the post-2006 rise in male firearm suicide represents a promising direction for future research.

% =============================================================================
% Back matter
% =============================================================================

\backmatter

\bibliographystyle{aea}
\bibliography{bibliography}

\end{document}
